\subsubsection{Décomposition harmonique du retour électronique}
\label{sec:el-espectro-retorno}

La trajectoire APP--TRIT--TPK produit d'abord le registre \(\tau_e\). La représentation harmonique agit ensuite sur ce registre et explicite l'information qu'évalue \(K_\Omega\). Dans cette réalisation, \(\pi\) et l'unité complexe sont les coordonnées déjà construites ; les caractères cycliques constituent un langage de lecture du mot produit.

Pour \(\tau\in\mathbb R^{12}\), avec des indices dans
\(\mathbb Z/12\mathbb Z\), nous fixons
\begin{equation}
 T_k(\tau)=\sum_{j=0}^{11}\tau_j
             \exp\!\left(-\frac{2\pi i kj}{12}\right),\qquad
 \theta_k=\frac{2\pi k}{12}.
 \label{eq:el-transformada-finita}
\end{equation}
La formule étend linéairement la représentation du registre
tritique à l'espace réel des mots. Les corrélations \(C_s\),
les fenêtres \(A_\ell\) et le lecteur \(P_6\) conservent les
définitions précédentes dans ce domaine élargi.

\begin{theorem}[Représentation spectrale du lecteur de retour]
\label{thm:el-lector-espectral}
Pour tout mot réel de longueur douze,
\begin{align}
 C_s(\tau)&=\frac1{12}\sum_{k=0}^{11}|T_k(\tau)|^2
                         \cos(s\theta_k),
 \label{eq:el-correlacion-espectral}\\
 \Omega(\tau)&=\frac1{12}\sum_{k=0}^{11}|T_k(\tau)|^2w_k,
 \label{eq:el-omega-espectral}\\
 w_k&=-2\cos\theta_k-4\cos(2\theta_k)-6\cos(3\theta_k),\notag\\
 P_6(\tau)&=\frac1{24}\sum_{k=0}^{11}|T_k(\tau)|^2(-1)^k.
 \label{eq:el-p6-espectral}
\end{align}
\end{theorem}
\begin{proof}
En développant le module au carré, on obtient
\[
 |T_k|^2=\sum_{a,b=0}^{11}\tau_a\tau_b
           \exp\!\left(-\frac{2\pi i k(a-b)}{12}\right).
\]
L'identité d'orthogonalité
\[
 \frac1{12}\sum_{k=0}^{11}
       \exp\!\left(\frac{2\pi i k n}{12}\right)
 =\begin{cases}1,&n\equiv0\pmod{12},\\0,&n\not\equiv0\pmod{12}
   \end{cases}
\]
se déduit en sommant une progression géométrique. En multipliant
l'expression de \(|T_k|^2\) par \(\exp(is\theta_k)\) et en sommant sur
\(k\), seuls subsistent les couples tels que \(a-b\equiv s\).
La somme résultante est \(C_s\) ; en prenant les parties réelles, on obtient
\eqref{eq:el-correlacion-espectral}. La combinaison
\(-2C_1-4C_2-6C_3\), déjà déduite des fenêtres, fournit
\eqref{eq:el-omega-espectral}. Enfin,
\(C_6=2P_6\) et \(\cos(6\theta_k)=(-1)^k\) donnent
\eqref{eq:el-p6-espectral}.
\end{proof}

\begin{corollary}[Trois modes du registre central]
\label{cor:el-tres-modos}
Le mot \(\tau_e=(1,1,1,-1,-1,-1)^2\) possède les puissances
spectrales
\begin{equation}
 |T_2|^2=64,\qquad |T_6|^2=16,\qquad |T_{10}|^2=64,
 \label{eq:el-potencias-espectrales}
\end{equation}
et les puissances restantes sont nulles. Son évaluation produit
\begin{equation}
 \begin{aligned}
 \Omega(\tau_e)&=\frac{64\cdot7+16\cdot4+64\cdot7}{12}=80,\\
 P_6(\tau_e)&=\frac{64+16+64}{24}=6.
 \end{aligned}
 \label{eq:el-evaluacion-espectral}
\end{equation}
\end{corollary}
\begin{proof}
Avec \(z_k=\exp(-i\theta_k)\), la répétition du sextuplet donne
\[
 T_k=(1+(-1)^k)(1-z_k^3)(1+z_k+z_k^2).
\]
Le premier facteur annule les indices impairs ; le second annule \(k=0,4,8\). Aux trois indices restants, on obtient
\[
 T_2=4-4i\sqrt3,\qquad T_6=4,\qquad T_{10}=4+4i\sqrt3.
\]
Leurs modules au carré sont ceux indiqués. La définition des
poids donne \(w_2=w_{10}=7\) et \(w_6=4\) ; tous ces indices sont
pairs. Les formules du théorème concluent l'évaluation.
\end{proof}

La substitution dans le caractère de retour fournit une deuxième expression exacte de sa dépendance à l'égard de l'histoire :
\begin{equation}
 \kappa_e=
 \frac1{12}\sum_{k=0}^{11}|T_k(\tau_e)|^2
       \left(w_k+\frac{\Delta_4}{2}(-1)^k\right)-54\alpha.
 \label{eq:el-kappa-espectral}
\end{equation}
Le terme \(-54\alpha\) conserve la provenance chronologique
de \(K_\Omega\). Les deux autres termes évaluent les puissances
harmoniques avec les poids des fenêtres et la corrélation
antipodale. La composition électronique complète
\eqref{eq:el-formula-completa} reçoit ainsi une lecture explicite
de ses modes de retour. La norme \(\sqrt3/4\), le transfert
\(\varphi/\pi^2\), les recensements \(22=27-5\) et \(15=5\cdot3\),
le défaut \(\Delta_4\) et le quotient \(R_{\mathrm{act}}\)
conservent les constructions précédemment démontrées et réunies dans
le tableau~\ref{tab:el-factores}.

\paragraph{Ordre, origine de phase et information conservée.}
Le mot alterné \(\tau_{\rm alt}=(1,-1)^6\) contient les mêmes six signes positifs et six signes négatifs que \(\tau_e\). Sa transformée n'a que \(T_6=12\), de sorte que \(\Omega(\tau_{\rm alt})=144\cdot4/12=48\). Ce témoin prouve la sensibilité à l'ordre d'un lecteur que l'histogramme des signes laisse indéterminé. Sa fonction est de délimiter le domaine du lecteur ; la sélection d'une espèce exige la section et la réalisation correspondantes.

Si \((S_a\tau)_j=\tau_{j+a}\), alors
\(T_k(S_a\tau)=\exp(ia\theta_k)T_k(\tau)\).
Les puissances et le caractère de retour conservent donc les
translations cycliques. Les coefficients complexes complets
reconstruisent le mot au moyen de
\[
 \tau_j=\frac1{12}\sum_{k=0}^{11}T_k(\tau)\exp(ij\theta_k),
\]
tandis que leurs modules au carré conservent l'autocorrélation.
La coordinatisation intégrale distingue en outre les historiques d'événements de
\eqref{eq:el-historias-centrales}. Cette hiérarchie précise quelle
information transmet chaque lecture du même transport.

Les identités précédentes possèdent des démonstrations algébriques sur
tout le domaine déclaré. Leur vérification reproductible évalue les
témoins dans \(\mathbb Z[z]/(z^4-z^2+1)\), vérifie les douze
translations et parcourt les \(4096\) mots binaires de longueur
douze pour contrôler l'identité entre fenêtres et corrélations.
Les trois modes sont des composantes du registre dodécaphasique ; les
réalisations d'états spatiaux liés utilisent en outre
les opérateurs dynamiques et d'interaction de leur propre domaine.
